\documentclass[12pt, a4paper]{article}

% ── Encodage & langue ──────────────────────────────────────────────────────────
\usepackage[utf8]{inputenc}
\usepackage[T1]{fontenc}
\usepackage[french]{babel}

% ── Mise en page ───────────────────────────────────────────────────────────────
\usepackage[top=2.5cm, bottom=2.5cm, left=2.5cm, right=2.5cm]{geometry}
\usepackage{parskip}
\usepackage{microtype}

% ── Tableaux & figures ─────────────────────────────────────────────────────────
\usepackage{booktabs}
\usepackage{tabularx}
\usepackage{multirow}
\usepackage{float}
\usepackage{graphicx}
\usepackage{array}

% ── Couleurs & code ────────────────────────────────────────────────────────────
\usepackage{xcolor}
\usepackage{listings}
\definecolor{codebg}{RGB}{245,245,245}
\definecolor{codegreen}{RGB}{0,128,0}
\definecolor{codegray}{RGB}{128,128,128}
\definecolor{accent}{RGB}{30,100,200}
\definecolor{highlight}{RGB}{255,243,220}
\definecolor{darkgreen}{RGB}{0,120,0}
\definecolor{darkred}{RGB}{180,0,0}

\lstset{
  backgroundcolor=\color{codebg},
  basicstyle=\ttfamily\small,
  commentstyle=\color{codegray},
  keywordstyle=\color{accent}\bfseries,
  stringstyle=\color{codegreen},
  frame=single,
  rulecolor=\color{gray!40},
  breaklines=true,
  captionpos=b,
  language=Python
}

% ── Boîtes colorées ────────────────────────────────────────────────────────────
\usepackage{tcolorbox}
\tcbuselibrary{skins}
\newtcolorbox{definitionbox}{
  colback=blue!5!white,
  colframe=accent,
  fonttitle=\bfseries,
  title=Définition,
  left=6pt, right=6pt
}
\newtcolorbox{resultbox}{
  colback=green!5!white,
  colframe=darkgreen,
  fonttitle=\bfseries,
  left=6pt, right=6pt
}

% ── Maths ──────────────────────────────────────────────────────────────────────
\usepackage{amsmath}
\usepackage{amssymb}

% ── Liens & méta ───────────────────────────────────────────────────────────────
\usepackage{hyperref}
\hypersetup{
  colorlinks=true,
  linkcolor=accent,
  urlcolor=accent,
  pdftitle={Prédiction de notes d'hôtels — Rapport},
  pdfauthor={},
}

% ── En-tête / pied de page ─────────────────────────────────────────────────────
\usepackage{fancyhdr}
\pagestyle{fancy}
\fancyhf{}
\rhead{\small Prédiction de notes d'hôtels}
\lhead{\small NLP — Rapport de projet}
\cfoot{\thepage}

% ══════════════════════════════════════════════════════════════════════════════
\begin{document}

% ── Page de titre ──────────────────────────────────────────────────────────────
\begin{titlepage}
  \centering
  \vspace*{2cm}
  {\Huge\bfseries Prédiction de notes d'hôtels\\[0.4em]
  à partir de commentaires textuels\par}
  \vspace{1cm}
  {\Large Rapport de projet — NLP\par}
  \vspace{2cm}
  \rule{\linewidth}{0.5pt}\\[0.5em]
  {\large\textit{Approches comparées :}\\[0.3em]
  TF-IDF \quad $\rightarrow$ \quad Sentence Transformers \quad $\rightarrow$ \quad RoBERTa fine-tuné\par}
  \rule{\linewidth}{0.5pt}
  \vfill
  {\large Mars 2026\par}
\end{titlepage}

\tableofcontents
\newpage

% ══════════════════════════════════════════════════════════════════════════════
\section{Présentation du projet}

\subsection{Objectif}

L'objectif de ce projet est de prédire automatiquement la note attribuée par un client
(de 1 à 5 étoiles) à partir du seul texte de son commentaire sur un hôtel, sans connaître
la note à l'avance.

Il s'agit d'un problème de \textbf{classification de texte supervisée} en \textbf{5 classes}
dans le domaine du \textbf{NLP} (Natural Language Processing — Traitement Automatique du Langage).
La classification est dite \textit{supervisée} car le modèle apprend à partir d'exemples
déjà étiquetés (texte + note connue), et \textit{multiclasse} car il y a 5 catégories
distinctes à prédire (et non 2 comme en classification binaire).

\subsection{Données}

Les données sont réparties en trois fichiers CSV, chacun ayant un rôle précis :

\begin{table}[H]
  \centering
  \begin{tabular}{llr}
    \toprule
    \textbf{Fichier} & \textbf{Usage} & \textbf{Exemples} \\
    \midrule
    \texttt{train\_hotel\_reviews.csv} & Entraîner les modèles & 18\,491 \\
    \texttt{valid\_hotel\_reviews.csv} & Évaluer pendant le développement & 1\,000 \\
    \texttt{test\_hotel\_reviews.csv}  & Évaluation finale unique & 1\,000 \\
    \bottomrule
  \end{tabular}
  \caption{Répartition des données}
\end{table}

Chaque ligne contient deux colonnes :
\begin{itemize}
  \item \texttt{Review} : texte libre du commentaire (en anglais)
  \item \texttt{Rating} : note entière de 1 à 5 étoiles
\end{itemize}

La règle fondamentale respectée tout au long du projet : le jeu de \textbf{test} n'est
consulté qu'une seule fois à la toute fin, après avoir figé le modèle définitif.
Toutes les décisions de développement (choix des hyperparamètres, comparaison des modèles)
sont prises sur le jeu de \textbf{validation}.

\subsection{Distribution des notes — Déséquilibre des classes}

\begin{table}[H]
  \centering
  \begin{tabular}{crc}
    \toprule
    \textbf{Note} & \textbf{Nb exemples (train)} & \textbf{Proportion} \\
    \midrule
    1\,$\star$ & 1\,310 &  7.1\,\% \\
    2\,$\star$ & 1\,630 &  8.8\,\% \\
    3\,$\star$ & 1\,982 & 10.7\,\% \\
    4\,$\star$ & 5\,465 & 29.6\,\% \\
    5\,$\star$ & 8\,104 & 43.8\,\% \\
    \midrule
    \textbf{Total} & \textbf{18\,491} & \textbf{100\,\%} \\
    \bottomrule
  \end{tabular}
  \caption{Distribution des notes dans le jeu d'entraînement}
\end{table}

Ce déséquilibre est le \textbf{défi principal} du projet : les notes 4\,$\star$ et
5\,$\star$ représentent à elles seules 73\,\% des données. Un modèle naïf qui prédirait
systématiquement 5\,$\star$ obtiendrait déjà 44\,\% d'accuracy (taux de bonnes prédictions)
sans rien apprendre. Les classes 2\,$\star$ et 3\,$\star$ sont structurellement
sous-représentées et donc difficiles à apprendre.

\subsection{Pipeline général du projet}

Voici la séquence complète depuis la donnée brute jusqu'à la prédiction finale :

\begin{center}
\ttfamily\small
\begin{tabular}{c}
\hline
Texte brut + note (CSV) \\
\hline
$\downarrow$ \\
\hline
Prétraitement du texte \\
(nettoyage ou texte brut selon l'approche) \\
\hline
$\downarrow$ \\
\hline
Représentation numérique du texte \\
(TF-IDF / Sentence Transformers / Tokenisation BERT) \\
\hline
$\downarrow$ \\
\hline
Modèle de classification \\
(LogReg / SVC / Réseau de neurones profond) \\
\hline
$\downarrow$ \\
\hline
Prédiction : note 1★ à 5★ \\
\hline
$\downarrow$ \\
\hline
Évaluation : Accuracy + F1 macro \\
\hline
\end{tabular}
\end{center}

\subsection{Métriques d'évaluation utilisées}

Deux métriques complémentaires sont calculées sur chaque modèle :

\begin{table}[H]
  \centering
  \begin{tabular}{lp{10cm}}
    \toprule
    \textbf{Métrique} & \textbf{Définition} \\
    \midrule
    \textbf{Accuracy} (taux de bonnes prédictions) &
    Proportion de notes exactement correctes parmi toutes les prédictions.
    $\text{Accuracy} = \frac{\text{nb bonnes prédictions}}{\text{nb total}}$.
    Peut être trompeuse si les classes sont déséquilibrées. \\[0.5em]
    \textbf{F1 macro} (moyenne harmonique précision/rappel, à poids égal par classe) &
    Pour chaque classe $c$, on calcule :
    $\text{F1}_c = \frac{2 \times P_c \times R_c}{P_c + R_c}$
    où $P_c$ (précision) = parmi les exemples prédits $c$, combien sont vraiment $c$ ?
    et $R_c$ (rappel) = parmi les vrais $c$, combien a-t-on correctement identifiés ?
    Le F1 macro est la moyenne de ces F1 sur les 5 classes — chaque classe compte
    autant, qu'elle ait 54 ou 484 exemples. \\
    \bottomrule
  \end{tabular}
  \caption{Métriques d'évaluation}
\end{table}

Le \textbf{F1 macro} est la métrique de référence car elle pénalise les modèles
qui ignorent les classes rares (2\,$\star$, 3\,$\star$), ce que l'accuracy seule ne détecterait pas.

\subsection{Plan de travail — Complexité croissante}

Trois approches sont comparées, classées par complexité croissante :

\begin{enumerate}
  \item \textbf{TF-IDF + modèles classiques} — représentation par comptage de mots, classifieur linéaire \hfill \textit{[réalisé]}
  \item \textbf{Sentence Transformers} — embeddings denses pré-entraînés, classifieur linéaire \hfill \textit{[réalisé]}
  \item \textbf{Fine-tuning RoBERTa} — transformer profond entraîné sur nos données \hfill \textit{[réalisé]}
\end{enumerate}

Cette progression permet de quantifier le gain apporté par chaque niveau de complexité
supplémentaire et de justifier l'usage d'architectures lourdes.

% ══════════════════════════════════════════════════════════════════════════════
\section{Étape 1 — Baseline TF-IDF + Modèles classiques}
\label{sec:tfidf}

\subsection{Principe et motivation}

TF-IDF (\textit{Term Frequency–Inverse Document Frequency} — Fréquence du terme pondérée
par l'inverse de sa fréquence dans les documents) est la représentation de texte classique
en NLP. Chaque commentaire est transformé en un vecteur numérique : chaque dimension
correspond à un mot du vocabulaire, et sa valeur mesure l'importance de ce mot dans ce document.

La formule est :
\begin{equation}
  \text{TF-IDF}(t, d) =
  \underbrace{\log(1 + f_{t,d})}_{\text{TF (fréquence du terme, avec lissage log)}}
  \times
  \underbrace{\log\frac{N}{|\{d : t \in d\}|}}_{\text{IDF (inverse de la fréquence documentaire)}}
\end{equation}

où $f_{t,d}$ est le nombre d'occurrences du terme $t$ dans le document $d$, et $N$ le nombre
total de documents. Un mot rare mais présent dans un avis (ex: ``immaculate'') obtient un score
élevé. Un mot très fréquent partout (ex: ``hotel'') obtient un score faible.

\textbf{Limites fondamentales de TF-IDF :}
\begin{itemize}
  \item \textbf{Pas de contexte} : ``not good'' et ``very good'' ont les mêmes dimensions
        actives, seuls les poids diffèrent — TF-IDF ne comprend pas la négation.
  \item \textbf{Pas de sémantique} : ``great'' et ``excellent'' sont deux dimensions
        entièrement distinctes, même s'ils signifient la même chose.
  \item \textbf{Représentation creuse} (sparse) : un vecteur de 20\,000 dimensions
        dont la grande majorité est à zéro (un commentaire n'utilise que quelques centaines
        de mots différents sur 20\,000 possibles).
\end{itemize}

\subsection{Prétraitement du texte}

Pour TF-IDF, un nettoyage \textbf{agressif} est appliqué avant la vectorisation.
Ce nettoyage ne sera \textbf{pas} appliqué aux approches Transformers (qui gèrent le texte brut).

\begin{lstlisting}[caption={Fonction de nettoyage pour TF-IDF}]
def clean_text(text):
    text = str(text)                      # Protection contre les valeurs NaN
    text = text.lower()                   # "Hotel" et "hotel" -> meme token
    text = re.sub(r'\d+', '', text)       # Supprimer les chiffres (peu informatifs)
    text = text.translate(str.maketrans( # Supprimer la ponctuation
        '', '', string.punctuation))
    text = re.sub(r'\s+', ' ', text).strip()  # Normaliser les espaces
    tokens = [w for w in text.split()
              if w not in STOPWORDS]      # Supprimer les stopwords
    return ' '.join(tokens)
\end{lstlisting}

\textbf{Justification de chaque étape :}
\begin{itemize}
  \item \textbf{Minuscules} : ``Hotel'', ``hotel'', ``HOTEL'' → même token, évite la redondance.
  \item \textbf{Suppression des chiffres} : les numéros de chambre, années, prix varient
        et n'apportent pas d'information sur la satisfaction.
  \item \textbf{Suppression des stopwords} (mots très fréquents sans valeur informative :
        ``the'', ``a'', ``is'', ``in'') : ces mots apparaissent dans presque tous les documents,
        leur IDF est quasi nul — ils n'aident pas à discriminer les classes.
  \item \textbf{Limite connue} : la suppression de ``not'' peut inverser le sens
        (``not good'' → ``good'') — c'est une limite acceptée pour la baseline.
\end{itemize}

\subsection{Vectorisation TF-IDF}

\begin{table}[H]
  \centering
  \begin{tabular}{lp{9cm}}
    \toprule
    \textbf{Paramètre} & \textbf{Valeur et justification} \\
    \midrule
    \texttt{max\_features=20000} &
      Garder les 20\,000 termes les plus fréquents. Réduit la dimensionnalité,
      supprime les mots ultra-rares (fautes de frappe, noms propres) qui
      ne généralisent pas. \\
    \texttt{ngram\_range=(1, 2)} &
      Unigrammes (``good'', ``clean'') ET bigrammes (``not good'', ``very clean'').
      Les bigrammes capturent les négations et expressions figées que les
      unigrammes seuls ne peuvent représenter. \\
    \texttt{sublinear\_tf=True} &
      Applique $\log(1 + \text{tf})$ — atténue l'effet des mots très répétés.
      ``clean'' mentionné 10 fois ne vaut pas 10$\times$ ``clean'' mentionné 1 fois. \\
    \texttt{min\_df=3} &
      Ignorer les mots présents dans moins de 3 documents. Supprime les hapax
      (mots uniques, fautes) qui ne généralisent pas. \\
    \bottomrule
  \end{tabular}
  \caption{Paramètres du TfidfVectorizer}
\end{table}

\textbf{Règle fondamentale — pas de fuite de données (data leakage) :}
Le vectoriseur apprend le vocabulaire et les scores IDF \textbf{uniquement sur train}.
Les jeux valid et test sont transformés avec ce vocabulaire figé.
Si on calculait les IDF sur valid ou test, le modèle ``verrait'' des données futures
et son évaluation serait artificiellement optimiste.

Résultat : matrices creuses de taille $(18\,491 \times 20\,000)$ pour train,
$(1\,000 \times 20\,000)$ pour valid et test.

\subsection{Pipeline TF-IDF complet}

\begin{center}
\ttfamily\small
\begin{tabular}{c}
Texte brut \\
$\downarrow$ \\
Nettoyage (lower + chiffres + ponctuation + stopwords) \\
$\downarrow$ \\
TfidfVectorizer.fit\_transform(train) \\
$\rightarrow$ Vecteur creux 20\,000 dimensions \\
$\downarrow$ \\
LogisticRegression / LinearSVC \\
$\downarrow$ \\
Prédiction : note 1 à 5 \\
\end{tabular}
\end{center}

\subsection{Modèles entraînés}

\subsubsection{Logistic Regression (régression logistique)}

La régression logistique est un modèle linéaire qui prédit la probabilité d'appartenance
à chaque classe. Pour 5 classes, elle entraîne 5 classificateurs ``one-vs-rest''
(un par classe, chacun contre toutes les autres). Le paramètre \texttt{C=1.0}
contrôle la régularisation L2 (pénalisation des poids trop grands pour éviter l'overfitting).

\subsubsection{LinearSVC (machine à vecteurs de support linéaire)}

LinearSVC cherche l'hyperplan qui sépare les classes avec la marge maximale.
C'est un classifieur souvent plus rapide que la régression logistique sur du texte
en haute dimension, mais il ne produit pas de probabilités — seulement des décisions.

\subsection{Résultats sur la validation}

\begin{table}[H]
  \centering
  \begin{tabular}{lcc}
    \toprule
    \textbf{Modèle} & \textbf{Accuracy (taux de bonnes prédictions)} & \textbf{F1 macro (moyenne F1 par classe)} \\
    \midrule
    \textbf{Logistic Regression} & \textbf{63.9\,\%} & \textbf{0.516} \\
    LinearSVC                    & 60.4\,\%          & 0.489          \\
    \bottomrule
  \end{tabular}
  \caption{Résultats sur le jeu de validation — TF-IDF}
\end{table}

La régression logistique surpasse LinearSVC malgré l'intuition contraire.
Sur des données déséquilibrées, LinearSVC tend à favoriser les grandes marges
qui profitent aux classes majoritaires, au détriment des classes rares.

\subsubsection{Détail par classe — Logistic Regression (validation)}

\begin{table}[H]
  \centering
  \begin{tabular}{lcccr}
    \toprule
    \textbf{Classe} & \textbf{Précision} & \textbf{Rappel} & \textbf{F1} & \textbf{Support} \\
    \midrule
    1\,$\star$ & 0.65 & 0.56 & 0.60 &  54 \\
    2\,$\star$ & 0.39 & 0.35 & 0.37 &  82 \\
    3\,$\star$ & 0.47 & 0.22 & \textcolor{darkred}{\textbf{0.29}} &  93 \\
    4\,$\star$ & 0.52 & 0.52 & 0.52 & 287 \\
    5\,$\star$ & 0.75 & 0.85 & \textcolor{darkgreen}{\textbf{0.79}} & 484 \\
    \midrule
    \textbf{Macro avg} & 0.56 & 0.50 & 0.52 & 1000 \\
    \bottomrule
  \end{tabular}
  \caption{Rapport de classification — Logistic Regression (validation)}
\end{table}

\textbf{Précision} (parmi les avis prédits $c$, combien sont vraiment $c$ ?) et
\textbf{rappel} (parmi les vrais $c$, combien a-t-on correctement identifiés ?) sont
les deux composantes du F1.

\subsection{Résultats sur le jeu de test}

\begin{table}[H]
  \centering
  \begin{tabular}{lcc}
    \toprule
    \textbf{Modèle} & \textbf{Accuracy} & \textbf{F1 macro} \\
    \midrule
    Logistic Regression & 61.3\,\% & 0.493 \\
    \bottomrule
  \end{tabular}
  \caption{Résultats finaux sur le jeu de test — TF-IDF}
\end{table}

\subsection{Analyse des erreurs}

Sur 1\,000 exemples de validation, 361 sont mal classés (36.1\,\%).
L'analyse des écarts montre que le modèle ne se trompe jamais grossièrement :

\begin{table}[H]
  \centering
  \begin{tabular}{cc}
    \toprule
    \textbf{Écart entre note réelle et prédite} & \textbf{Nb erreurs} \\
    \midrule
    1\,$\star$ (erreur légère) & 304 (84\,\%) \\
    2\,$\star$ (erreur modérée) & 44 (12\,\%) \\
    $\geq$3\,$\star$ (erreur grave) & 13 (4\,\%) \\
    \bottomrule
  \end{tabular}
  \caption{Répartition des erreurs par écart — Logistic Regression}
\end{table}

\subsection{Observations et limites}

\begin{itemize}
  \item La note 3\,$\star$ (F1 = 0.29, rappel = 22\,\%) est la plus difficile : un avis
        mitigé contient des mots positifs ET négatifs, TF-IDF ne comprend pas leur équilibre.
  \item La note 5\,$\star$ est la mieux prédite (F1 = 0.79) grâce à sa fréquence et
        à un vocabulaire très distinctif (``perfect'', ``excellent'', ``amazing'').
  \item Le déséquilibre des classes est le principal facteur limitant :
        le modèle tire ses prédictions vers les classes fréquentes.
  \item \textbf{Ce score sert de baseline} (point de référence) : toute approche
        ultérieure doit dépasser 61.3\,\% d'accuracy et 0.493 de F1 macro pour être utile.
\end{itemize}

% ══════════════════════════════════════════════════════════════════════════════
\section{Étape 2 — Sentence Transformers}
\label{sec:sentencetransformers}

\subsection{Principe et motivation}

Sentence Transformers est une famille de modèles basés sur BERT
(Bidirectional Encoder Representations from Transformers — représentations
bidirectionnelles de transformers encodeurs) ré-entraînés pour produire
un \textbf{vecteur dense unique} représentant le sens global d'une phrase entière.

Contrairement à TF-IDF qui compte les mots, Sentence Transformers
\textbf{comprend le sens} grâce à un mécanisme d'attention bidirectionnelle :
chaque mot est représenté en tenant compte de tous les autres mots de la phrase.

\textbf{Modèle utilisé :} \texttt{all-MiniLM-L6-v2}
\begin{itemize}
  \item 6 couches Transformer (version allégée de BERT à 12 couches)
  \item Vecteurs de 384 dimensions (denses — toutes les valeurs sont significatives)
  \item Taille : $\approx$80\,Mo — exécutable sur CPU sans GPU
  \item Pré-entraîné par apprentissage contrastif sur plus d'un milliard de paires de phrases :
        deux phrases de sens proche $\rightarrow$ vecteurs proches dans l'espace,
        deux phrases de sens opposé $\rightarrow$ vecteurs éloignés.
\end{itemize}

\begin{table}[H]
  \centering
  \begin{tabular}{lll}
    \toprule
    \textbf{Critère} & \textbf{TF-IDF} & \textbf{Sentence Transformers} \\
    \midrule
    Représentation    & Vecteur creux (20\,000 dim) & Vecteur dense (384 dim) \\
    Contexte          & Aucun & Bidirectionnel \\
    Sémantique        & Non (\og great\fg{} $\neq$ \og excellent\fg{}) & Oui \\
    \og not good\fg{} & Similaire à \og very good\fg{} & Vecteur opposé \\
    GPU requis        & Non & Non \\
    Fine-tuning tâche & Non & Non \\
    \bottomrule
  \end{tabular}
  \caption{Comparaison TF-IDF vs Sentence Transformers}
\end{table}

\subsection{Prétraitement — Texte brut}

Contrairement à TF-IDF, \textbf{aucun nettoyage agressif} n'est appliqué ici.
Le modèle a été pré-entraîné sur du texte brut et a appris que :
\begin{itemize}
  \item ``not good'' $\neq$ ``good'' (la négation change le sens — conserver ``not'')
  \item ``Very clean!'' $\neq$ ``clean'' (l'intensité et la ponctuation portent de l'information)
\end{itemize}
Supprimer les stopwords ou la ponctuation \textbf{dégraderait} les embeddings.
Seule la gestion des valeurs nulles (NaN) et un \texttt{strip()} sont effectués.

\subsection{Pipeline Sentence Transformers complet}

\begin{center}
\ttfamily\small
\begin{tabular}{c}
Texte brut (strip uniquement) \\
$\downarrow$ \\
all-MiniLM-L6-v2 (modèle pré-entraîné, poids figés) \\
$\downarrow$ \\
Vecteur dense de 384 dimensions par avis \\
(sauvegardé en .npy pour éviter le recalcul) \\
$\downarrow$ \\
LogisticRegression (C=5.0) \\
$\downarrow$ \\
Prédiction : note 1 à 5 \\
\end{tabular}
\end{center}

\textbf{Pourquoi sauvegarder les embeddings en .npy ?}
Encoder 18\,491 textes avec un transformer prend plusieurs minutes même sur CPU.
Le format \texttt{.npy} (numpy binaire) permet de charger instantanément les vecteurs
lors des exécutions suivantes du notebook.

\textbf{Pourquoi C=5.0 au lieu de 1.0 ?}
Les embeddings (384 dimensions denses) sont beaucoup moins bruités que les vecteurs
TF-IDF (20\,000 dimensions creuses). La régularisation peut donc être réduite
(C plus grand = moins de contrainte) car les features sont déjà bien structurées.

\subsection{Résultats sur la validation}

\begin{table}[H]
  \centering
  \begin{tabular}{lcc}
    \toprule
    \textbf{Modèle} & \textbf{Accuracy} & \textbf{F1 macro} \\
    \midrule
    Sentence Transformers + LogReg & 63.1\,\% & 0.510 \\
    \bottomrule
  \end{tabular}
  \caption{Résultats sur le jeu de validation — Sentence Transformers}
\end{table}

\subsubsection{Détail par classe (validation)}

\begin{table}[H]
  \centering
  \begin{tabular}{lcccr}
    \toprule
    \textbf{Classe} & \textbf{Précision} & \textbf{Rappel} & \textbf{F1} & \textbf{Support} \\
    \midrule
    1\,$\star$ & 0.70 & 0.57 & 0.63 &  54 \\
    2\,$\star$ & 0.43 & 0.51 & 0.47 &  82 \\
    3\,$\star$ & 0.27 & 0.11 & \textcolor{darkred}{\textbf{0.15}} &  93 \\
    4\,$\star$ & 0.54 & 0.49 & 0.52 & 287 \\
    5\,$\star$ & 0.73 & 0.84 & \textcolor{darkgreen}{\textbf{0.78}} & 484 \\
    \midrule
    \textbf{Macro avg} & 0.53 & 0.51 & 0.51 & 1000 \\
    \bottomrule
  \end{tabular}
  \caption{Rapport de classification — Sentence Transformers (validation)}
\end{table}

\subsection{Résultats sur le jeu de test}

\begin{table}[H]
  \centering
  \begin{tabular}{lcccr}
    \toprule
    \textbf{Classe} & \textbf{Précision} & \textbf{Rappel} & \textbf{F1} & \textbf{Support} \\
    \midrule
    1\,$\star$ & 0.67 & 0.67 & 0.67 &  57 \\
    2\,$\star$ & 0.51 & 0.53 & 0.52 &  81 \\
    3\,$\star$ & 0.53 & 0.28 & \textcolor{darkred}{\textbf{0.37}} & 109 \\
    4\,$\star$ & 0.51 & 0.50 & 0.51 & 287 \\
    5\,$\star$ & 0.73 & 0.82 & \textcolor{darkgreen}{\textbf{0.77}} & 466 \\
    \midrule
    \textbf{Macro avg} & 0.59 & 0.56 & 0.57 & 1000 \\
    \midrule
    \textbf{Accuracy globale} & \multicolumn{4}{c}{\textbf{63.5\,\%} — F1 macro : \textbf{0.566}} \\
    \bottomrule
  \end{tabular}
  \caption{Rapport de classification — Sentence Transformers (test)}
\end{table}

\subsection{Analyse comparative TF-IDF vs Sentence Transformers}

\begin{table}[H]
  \centering
  \begin{tabular}{lccc}
    \toprule
    \textbf{Classe} & \textbf{TF-IDF F1} & \textbf{SentTrans F1} & \textbf{Différence} \\
    \midrule
    1\,$\star$ & 0.60 & 0.63 & \textcolor{darkgreen}{+0.03} \\
    2\,$\star$ & 0.37 & 0.47 & \textcolor{darkgreen}{+0.10} \\
    3\,$\star$ & 0.29 & 0.15 & \textcolor{darkred}{$-$0.14} \\
    4\,$\star$ & 0.52 & 0.52 & 0.00 \\
    5\,$\star$ & 0.79 & 0.78 & \textcolor{darkred}{$-$0.01} \\
    \bottomrule
  \end{tabular}
  \caption{Comparaison F1 par classe : TF-IDF vs Sentence Transformers (validation)}
\end{table}

\subsection{Observations}

\begin{itemize}
  \item Sur la \textbf{validation}, Sentence Transformers est légèrement inférieur à
        TF-IDF (F1 macro : 0.510 vs 0.516). C'est un résultat paradoxal.
  \item Sur le \textbf{test}, Sentence Transformers surpasse nettement TF-IDF
        (F1 macro : \textbf{0.566} vs 0.493, soit +0.073), ce qui indique
        une \textbf{meilleure généralisation} — les embeddings sémantiques capturent
        des patterns plus robustes que les fréquences de mots.
  \item La note 2\,$\star$ bénéficie le plus de l'approche sémantique (+0.10 de F1) :
        le modèle comprend le contexte négatif même quand les mots négatifs sont
        accompagnés de nuances.
  \item La note 3\,$\star$ s'effondre (F1 = 0.15 sur valid, rappel = 11\,\%) : le modèle
        évite de prédire cette classe ambiguë et préfère systématiquement 2\,$\star$ ou 4\,$\star$.
        Cela s'explique par l'absence de fine-tuning : \texttt{all-MiniLM-L6-v2} est
        optimisé pour la similarité de phrases, pas pour la classification de sentiments graduée.
  \item \textbf{Limite fondamentale} : les deux approches (TF-IDF et Sentence Transformers)
        utilisent le modèle \textbf{sans l'adapter} à notre tâche. Le fine-tuning
        de la section suivante va remédier à cela.
\end{itemize}

% ══════════════════════════════════════════════════════════════════════════════
\section{Étape 3 — Fine-tuning RoBERTa}
\label{sec:bert}

\subsection{Principe — Qu'est-ce que le fine-tuning ?}

Le fine-tuning (affinage) consiste à prendre un modèle pré-entraîné sur de très grandes
quantités de données génériques, puis à le \textbf{ré-entraîner entièrement} sur nos
données spécifiques. Tous les poids du modèle sont mis à jour — pas seulement une tête
de classification ajoutée par-dessus.

\textbf{Modèle utilisé : RoBERTa-base}
(Robustly Optimized BERT Pretraining Approach — approche d'entraînement BERT robustement optimisée)
\begin{itemize}
  \item 125 millions de paramètres
  \item 12 couches Transformer (encodeur bidirectionnel)
  \item Vocabulaire BPE (Byte-Pair Encoding — encodage par paires d'octets) de 50\,265 tokens
  \item Pré-entraîné sur 160\,Go de texte (Wikipedia, livres, pages web)
  \item Amélioration de BERT : plus de données, meilleure stratégie de masquage,
        suppression des segments de phrases (plus simple et plus efficace)
\end{itemize}

\subsection{Architecture détaillée de RoBERTa pour la classification}

\begin{center}
\ttfamily\small
\begin{tabular}{c}
Texte brut \\
$\downarrow$ \\
Tokenizer BPE (Byte-Pair Encoding) \\
``unbelievable'' -> [``un'', ``believ'', ``able''] \\
$\downarrow$ \\
input\_ids [CLS, t1, t2, ..., tN, SEP, PAD, PAD, ...] + attention\_mask \\
(max\_length = 256 tokens) \\
$\downarrow$ \\
Embeddings (token + position) \\
$\downarrow$ \\
12 couches Transformer Encoder \\
Chaque couche : Multi-Head Self-Attention + Feed-Forward + LayerNorm \\
$\downarrow$ \\
Vecteur du token [CLS] (768 dimensions) \\
[CLS] = Classification token, token special en debut de sequence \\
      = resume le sens global de toute la phrase \\
$\downarrow$ \\
Dense(768 -> 768) + Dropout + Dense(768 -> 5) \\
= Tete de classification (initialisee aleatoirement, apprise par fine-tuning) \\
$\downarrow$ \\
Logits (5 scores bruts, un par classe) \\
$\downarrow$ \\
Softmax -> Probabilites -> argmax -> Classe predite (0 a 4) \\
$\downarrow$ \\
+1 -> Note predite (1 a 5) \\
\end{tabular}
\end{center}

\subsubsection{Le token [CLS] — Classification Token}

Le token \texttt{[CLS]} (Classification Token — token de classification) est un token
spécial inséré automatiquement \textbf{au début} de chaque séquence par le tokenizer.
Après avoir traversé les 12 couches Transformer, ce token a été mis à jour par
l'attention de tous les autres tokens de la séquence. Son vecteur final (768 dimensions)
encode donc une représentation globale de toute la phrase, utilisée comme entrée
de la tête de classification.

\subsubsection{Multi-Head Self-Attention — Mécanisme d'attention}

L'attention (Attention) est le mécanisme central des Transformers. Dans chaque couche,
chaque token ``regarde'' tous les autres tokens de la séquence et décide combien
de poids leur accorder pour construire sa propre représentation.

``Multi-head'' (multi-têtes) signifie que ce mécanisme est appliqué en parallèle
12 fois avec des projections différentes, permettant de capturer différents types
de relations entre les mots (syntaxe, sémantique, coréférence...).

\subsection{Prétraitement — Tokenisation}

Contrairement à TF-IDF, \textbf{aucun nettoyage manuel} n'est appliqué.
Le tokenizer RoBERTa gère lui-même la normalisation.

\begin{lstlisting}[caption={Tokenisation pour RoBERTa}]
tokenizer = RobertaTokenizer.from_pretrained('roberta-base')

def tokenize(texts):
    return tokenizer(
        texts,
        max_length=256,        # Tronquer a 256 tokens (limite memoire GPU)
        truncation=True,       # Couper les textes trop longs
        padding='max_length',  # Padder les textes courts a 256
        return_tensors='pt'    # Retourner des tenseurs PyTorch
    )
\end{lstlisting}

\textbf{Pourquoi MAX\_LEN=256 et pas 512 ?}
La mémoire GPU (VRAM) consommée par l'attention est quadratique avec la longueur
de la séquence ($O(n^2)$) : doubler la longueur quadruple la mémoire.
MAX\_LEN=256 permet un batch de 32 exemples sur le GPU T4 (15.6\,Go VRAM),
contre un batch de 8 seulement avec MAX\_LEN=512.
La grande majorité des avis hôtels étant inférieurs à 256 tokens, la perte
d'information par troncature est négligeable.

\textbf{Décalage des labels :}
PyTorch \texttt{CrossEntropyLoss} (perte d'entropie croisée) exige des labels
entiers de $0$ à $N-1$. Les notes 1--5 sont donc converties en labels 0--4
($\text{label} = \text{Rating} - 1$). À l'évaluation, on reconvertit : $\text{note} = \text{pred} + 1$.

\subsection{Dataset PyTorch et DataLoader}

Pour que PyTorch puisse itérer sur les données par mini-batches (petits lots),
les données tokenisées sont encapsulées dans une classe \texttt{Dataset} :

\begin{lstlisting}[caption={Dataset PyTorch pour les avis d'hôtels}]
class HotelDataset(Dataset):
    def __init__(self, encodings, labels):
        self.input_ids      = encodings['input_ids']
        self.attention_mask = encodings['attention_mask']
        self.labels = torch.tensor(labels, dtype=torch.long)

    def __len__(self): return len(self.labels)

    def __getitem__(self, idx):
        return {
            'input_ids'      : self.input_ids[idx],
            'attention_mask' : self.attention_mask[idx],
            'labels'         : self.labels[idx]
        }
\end{lstlisting}

Le \texttt{DataLoader} (chargeur de données) utilise cette classe pour former des
batches de taille 32, en mélangeant les exemples à chaque époque (\texttt{shuffle=True}
sur train uniquement).

\textbf{Pourquoi batch\_size=32 ?}
Plus le batch est grand, plus les gradients sont stables (estimés sur plus d'exemples),
et plus l'entraînement est rapide (meilleure utilisation parallèle du GPU).
32 est la valeur maximale qui tient en VRAM avec RoBERTa-base et MAX\_LEN=256 sur T4.

\subsection{Optimiseur, scheduler et hyperparamètres}

\begin{table}[H]
  \centering
  \begin{tabular}{lp{9cm}}
    \toprule
    \textbf{Hyperparamètre} & \textbf{Valeur et justification} \\
    \midrule
    \textbf{Optimiseur} & AdamW (Adam avec Weight Decay découplé — régularisation L2).
      Standard pour le fine-tuning de transformers. \\
    \textbf{Learning rate} & $2 \times 10^{-5}$.
      Très petit pour préserver les représentations pré-entraînées.
      Un LR trop grand causerait le \textit{catastrophic forgetting}
      (oubli catastrophique des connaissances acquises). \\
    \textbf{Weight decay} & 0.01. Pénalise les poids trop grands → réduit l'overfitting. \\
    \textbf{Époques} & 3. L'époque 2 donne le meilleur F1 valid ;
      l'époque 3 montre un léger overfitting (valid loss qui remonte). \\
    \textbf{Warmup} & 10\,\% des steps ($\approx$173 steps sur 1\,734 totaux).
      Le LR monte de 0 à $2 \times 10^{-5}$ linéairement, puis redescend vers 0.
      Stabilise les premières étapes où la tête de classification est aléatoire. \\
    \textbf{Gradient clipping} & Norme maximale = 1.0.
      Plafonne les gradients pour éviter les explosions de gradients,
      phénomène courant dans les transformers profonds. \\
    \bottomrule
  \end{tabular}
  \caption{Hyperparamètres RoBERTa-base}
\end{table}

\subsection{Boucle d'entraînement}

À chaque époque, la boucle d'entraînement suit ce schéma sur chaque batch :

\begin{center}
\ttfamily\small
\begin{tabular}{c}
optimizer.zero\_grad() \\
(reinitialiser les gradients accumules du batch precedent) \\
$\downarrow$ \\
Forward pass : model(input\_ids, attention\_mask, labels) \\
-> Calcul des logits (5 scores) + CrossEntropyLoss \\
(la loss mesure l'ecart entre la prediction et la vraie note) \\
$\downarrow$ \\
loss.backward() \\
(retropropagation : calcul du gradient de chaque parametre) \\
$\downarrow$ \\
clip\_grad\_norm\_(model.parameters(), 1.0) \\
(ecreter les gradients trop grands) \\
$\downarrow$ \\
optimizer.step() \\
(mettre a jour les 125M parametres dans la direction opposee au gradient) \\
$\downarrow$ \\
scheduler.step() \\
(ajuster le learning rate selon le profil warmup + descente lineaire) \\
\end{tabular}
\end{center}

Après chaque époque, le modèle est évalué sur le jeu de validation.
Le meilleur checkpoint (selon le F1 macro valid) est sauvegardé dans \texttt{best\_roberta.pt}.

\textbf{Pourquoi sauvegarder selon le F1 et non l'accuracy ?}
Sur des données déséquilibrées, l'accuracy peut augmenter même si les classes rares
se dégradent (car les prédictions correctes sur 5\,$\star$ compensent).
Le F1 macro pénalise les mauvaises performances sur les classes minoritaires.

\subsection{Courbes d'apprentissage}

\begin{table}[H]
  \centering
  \begin{tabular}{ccccc}
    \toprule
    \textbf{Époque} & \textbf{Train loss} & \textbf{Valid loss} & \textbf{Accuracy valid} & \textbf{F1 macro valid} \\
    \midrule
    1 & 0.9442 & 0.7918 & 65.8\,\% & 0.6119 \\
    \textbf{2} & \textbf{0.6955} & \textbf{0.7038} & \textbf{69.4\,\%} & \textbf{0.6316} $\leftarrow$ \textit{sauvegardé} \\
    3 & 0.6064 & 0.7205 & 69.9\,\% & 0.6295 \\
    \bottomrule
  \end{tabular}
  \caption{Courbes d'apprentissage — RoBERTa-base (3 époques)}
\end{table}

\textbf{Interprétation :}
\begin{itemize}
  \item \textbf{Époque 1} : le modèle apprend rapidement les patterns de base.
        La train loss (0.9442) est encore élevée — le modèle fait encore beaucoup d'erreurs.
  \item \textbf{Époque 2} : meilleur équilibre. La valid loss (0.7038) est à son minimum
        et le F1 macro valid (0.6316) est au plus haut → checkpoint sauvegardé.
  \item \textbf{Époque 3} : la train loss continue de baisser (0.6064) mais la valid loss
        remonte légèrement (0.7205) et le F1 valid baisse (0.6295).
        C'est le signe d'un début d'\textbf{overfitting} (sur-apprentissage) :
        le modèle mémorise les données d'entraînement au détriment de la généralisation.
\end{itemize}

\subsection{Résultats finaux — Meilleur checkpoint (époque 2)}

\subsubsection{Validation}

\begin{table}[H]
  \centering
  \begin{tabular}{lcccr}
    \toprule
    \textbf{Classe} & \textbf{Précision} & \textbf{Rappel} & \textbf{F1} & \textbf{Support} \\
    \midrule
    1\,$\star$ & 0.80 & 0.69 & 0.74 &  54 \\
    2\,$\star$ & 0.52 & 0.62 & 0.56 &  82 \\
    3\,$\star$ & 0.55 & 0.41 & \textcolor{darkred}{\textbf{0.47}} &  93 \\
    4\,$\star$ & 0.59 & 0.54 & 0.57 & 287 \\
    5\,$\star$ & 0.79 & 0.85 & \textcolor{darkgreen}{\textbf{0.82}} & 484 \\
    \midrule
    \textbf{Macro avg} & 0.65 & 0.62 & 0.63 & 1000 \\
    \midrule
    \textbf{Accuracy} & \multicolumn{4}{c}{\textbf{69.4\,\%} — F1 macro : \textbf{0.6316}} \\
    \bottomrule
  \end{tabular}
  \caption{Rapport de classification — RoBERTa-base (validation)}
\end{table}

\subsubsection{Test}

\begin{table}[H]
  \centering
  \begin{tabular}{lcccr}
    \toprule
    \textbf{Classe} & \textbf{Précision} & \textbf{Rappel} & \textbf{F1} & \textbf{Support} \\
    \midrule
    1\,$\star$ & 0.69 & 0.75 & 0.72 &  57 \\
    2\,$\star$ & 0.56 & 0.60 & 0.58 &  81 \\
    3\,$\star$ & 0.61 & 0.35 & \textcolor{darkred}{\textbf{0.44}} & 109 \\
    4\,$\star$ & 0.60 & 0.60 & 0.60 & 287 \\
    5\,$\star$ & 0.81 & 0.86 & \textcolor{darkgreen}{\textbf{0.83}} & 466 \\
    \midrule
    \textbf{Macro avg} & 0.65 & 0.63 & 0.64 & 1000 \\
    \midrule
    \textbf{Accuracy} & \multicolumn{4}{c}{\textbf{70.4\,\%} — F1 macro : \textbf{0.636}} \\
    \bottomrule
  \end{tabular}
  \caption{Rapport de classification — RoBERTa-base (test) ← résultat final}
\end{table}

\subsection{Analyse des erreurs}

Sur 1\,000 exemples de validation, 306 sont mal classés (30.6\,\%).
94\,\% des erreurs ont un écart d'une seule étoile entre note réelle et note prédite —
le modèle ne se trompe jamais grossièrement, uniquement sur les frontières floues
(avis mitigés difficiles à distinguer entre deux notes adjacentes).

\textbf{Confusions principales observées :}
\begin{itemize}
  \item Réel 3\,$\star$ $\rightarrow$ prédit 4\,$\star$ : erreur la plus fréquente.
        Un avis globalement satisfait avec quelques réserves peut sembler positif.
  \item Réel 4\,$\star$ $\rightarrow$ prédit 5\,$\star$ : frontière subjective entre
        ``très bien'' et ``parfait''.
  \item Réel 3\,$\star$ : le rappel de 35\,\% signifie que 65 vrais avis mitigés
        sur 100 sont mal classés — c'est le point faible structurel du modèle.
\end{itemize}

\subsection{Observations}

\begin{itemize}
  \item \textbf{Gain majeur} par rapport aux approches précédentes : +9.1 points d'accuracy
        et +0.143 de F1 macro par rapport à TF-IDF.
  \item \textbf{Le fine-tuning fait la différence} : Sentence Transformers (sans fine-tuning)
        plafonnait à 63.5\,\%. RoBERTa fine-tuné atteint 70.4\,\% — le même type de modèle
        mais adapté à la tâche spécifique.
  \item \textbf{La note 3\,$\star$ reste le point faible} (F1 = 0.44, rappel = 35\,\%) :
        les avis ambigus sont intrinsèquement difficiles à classer, même pour un humain.
  \item L'overfitting apparaît dès l'époque 3 : 3 époques est le bon compromis
        pour ce dataset de taille modérée (18\,491 exemples).
\end{itemize}

% ══════════════════════════════════════════════════════════════════════════════
\section{Comparaison générale des approches}
\label{sec:comparaison}

\subsection{Tableau récapitulatif complet}

\begin{table}[H]
  \centering
  \begin{tabularx}{\textwidth}{Xcccc}
    \toprule
    \textbf{Approche} & \textbf{Acc. valid} & \textbf{F1 valid} & \textbf{Acc. test} & \textbf{F1 test} \\
    \midrule
    TF-IDF + LogReg              & 63.9\,\% & 0.516 & 61.3\,\% & 0.493 \\
    TF-IDF + LinearSVC           & 60.4\,\% & 0.489 & —        & — \\
    Sentence Transformers + LogReg & 63.1\,\% & 0.510 & 63.5\,\% & 0.566 \\
    \textbf{RoBERTa-base fine-tuné} & \textbf{69.4\,\%} & \textbf{0.632} & \textbf{70.4\,\%} & \textbf{0.636} \\
    \bottomrule
  \end{tabularx}
  \caption{Tableau comparatif de toutes les approches}
\end{table}

\subsection{Progression du F1 macro par classe}

\begin{table}[H]
  \centering
  \begin{tabular}{lcccc}
    \toprule
    \textbf{Classe} & \textbf{TF-IDF} & \textbf{SentTrans} & \textbf{RoBERTa} & \textbf{Progression} \\
    \midrule
    1\,$\star$ & 0.62 & 0.67 & 0.72 & \textcolor{darkgreen}{+0.10} \\
    2\,$\star$ & 0.39 & 0.52 & 0.58 & \textcolor{darkgreen}{+0.19} \\
    3\,$\star$ & 0.18 & 0.37 & 0.44 & \textcolor{darkgreen}{+0.26} \\
    4\,$\star$ & 0.49 & 0.51 & 0.60 & \textcolor{darkgreen}{+0.11} \\
    5\,$\star$ & 0.78 & 0.77 & 0.83 & \textcolor{darkgreen}{+0.05} \\
    \bottomrule
  \end{tabular}
  \caption{Évolution du F1 par classe (test) — de TF-IDF à RoBERTa}
\end{table}

\subsection{Analyse de la progression}

\subsubsection{Deux paliers de performance distincts}

Les résultats révèlent deux paliers nets :

\begin{itemize}
  \item \textbf{Palier 1 — Méthodes sans fine-tuning} (61\,\%--63\,\%) :
        TF-IDF et Sentence Transformers plafonnent dans cette zone malgré des approches
        fondamentalement différentes. La limite vient de l'absence d'adaptation à la tâche
        spécifique de classification de notes d'hôtels.
  \item \textbf{Palier 2 — Fine-tuning} (70\,\%) :
        RoBERTa fine-tuné franchit un seuil que les méthodes sans fine-tuning ne peuvent
        pas atteindre. L'entraînement sur nos données spécifiques permet au modèle
        de capturer des patterns propres aux avis hôteliers.
\end{itemize}

\subsubsection{La représentation est le facteur dominant}

Le gain le plus important provient non pas du changement de classifieur
(LogReg vs SVC : +3 points), mais du changement de représentation :
\begin{itemize}
  \item TF-IDF $\rightarrow$ Sentence Transformers : +2.2 points d'accuracy sur test
  \item Sentence Transformers $\rightarrow$ RoBERTa fine-tuné : +6.9 points d'accuracy sur test
\end{itemize}

\subsubsection{Le déséquilibre des classes est le problème non résolu}

La note 3\,$\star$ reste le point faible de \textbf{toutes} les approches,
progressant de F1=0.18 (TF-IDF) à F1=0.44 (RoBERTa). Bien que la progression
soit significative (+0.26), le rappel de 35\,\% pour 3\,$\star$ indique que
65\,\% des vrais avis mitigés sont encore mal classés. Les class weights
(pondération des classes pendant la perte) n'ont pas été appliqués
à RoBERTa dans cette version.

\subsection{Comparaison par type d'approche}

\begin{table}[H]
  \centering
  \begin{tabular}{lllll}
    \toprule
    \textbf{Critère} & \textbf{TF-IDF} & \textbf{SentTrans} & \textbf{RoBERTa} \\
    \midrule
    Représentation     & Creuse, 20k dim & Dense, 384 dim  & Dense, 768 dim \\
    Contextuelle       & Non             & Oui             & Oui \\
    Fine-tuning tâche  & Non             & Non             & Oui \\
    GPU requis         & Non             & Non             & Recommandé \\
    Temps d'entraînement & Secondes      & Minutes (CPU)   & $\approx$45 min (GPU T4) \\
    Accuracy test      & 61.3\,\%        & 63.5\,\%        & \textbf{70.4\,\%} \\
    F1 macro test      & 0.493           & 0.566           & \textbf{0.636} \\
    \bottomrule
  \end{tabular}
  \caption{Comparaison qualitative des approches}
\end{table}

% ══════════════════════════════════════════════════════════════════════════════
\section*{Conclusion}
\addcontentsline{toc}{section}{Conclusion}

Ce projet compare trois approches de NLP pour la classification d'avis hôteliers
en 5 classes de sentiment (1 à 5 étoiles), sur un dataset de 18\,491 exemples
d'entraînement fortement déséquilibré (44\,\% de notes 5\,$\star$).

\textbf{Résultats clés :}
\begin{enumerate}
  \item \textbf{RoBERTa-base fine-tuné} est le meilleur modèle avec 70.4\,\%
        d'accuracy et un F1 macro (moyenne harmonique précision/rappel à poids égal
        par classe) de 0.636 sur le jeu de test.
  \item \textbf{La représentation est le facteur dominant} : le passage de TF-IDF
        à RoBERTa apporte +9.1 points d'accuracy, bien plus que tout changement
        de classifieur classique.
  \item \textbf{Le fine-tuning est essentiel} : Sentence Transformers sans fine-tuning
        (63.5\,\%) plafonne nettement en dessous de RoBERTa fine-tuné (70.4\,\%),
        malgré une architecture similaire.
  \item \textbf{La note 3\,$\star$ reste le défi central} : avec un rappel de 35\,\%,
        les avis ambigus sont intrinsèquement difficiles à classer, même pour un modèle
        de 125 millions de paramètres.
\end{enumerate}

\textbf{Pistes d'amélioration identifiées :}
\begin{itemize}
  \item Fusionner les classes 3\,$\star$ et 4\,$\star$ pour passer à 4 classes distinctes.
        La note 3\,$\star$ est structurellement ambiguë : ses avis partagent un vocabulaire
        proche de la note 4\,$\star$ (satisfaction mitigée mais globalement positive),
        ce qui explique qu'elle soit la classe la plus mal prédite dans tous les modèles testés.
        Cette fusion réduirait mécaniquement les confusions et améliorerait le F1 macro.
  \item Appliquer des class weights (pondération des classes minoritaires) dans la
        fonction de perte CrossEntropyLoss pour RoBERTa, afin de pénaliser davantage
        les erreurs sur les notes rares.
  \item Combinaison par ensemble (vote ou moyenne des logits) des prédictions de
        plusieurs modèles pour gagner quelques points supplémentaires.
\end{itemize}

\end{document}
